\chapter{The Finite-Element Function Spaces}
\label{ch:functionspaces}

\chapterorient{The previous chapters turned each Maxwell problem into a weak form
$a(u,v)=\ell(v)$ and a Galerkin recipe: pick a finite-dimensional space $V_h$,
expand the unknown in a basis, and solve a linear system. We left one question
open---\emph{which} space? This chapter answers it. There are four canonical
finite-element spaces, one for each kind of physical field, and the choice is not
a matter of taste: a field that must be continuous in its tangential part lives in
a different space from one that must be continuous in its normal part. Put a field
in the wrong space and the simulator returns nonsense---spurious modes, locked
solutions, fields that cannot represent a constant. Get it right and the
discretization inherits the exact structure of Maxwell's equations. By the end you
will know why electrostatics uses a scalar nodal space while every magnetic and
wave problem uses edge elements, you will be able to count the degrees of freedom
of each space on a mesh by hand, and you will have met the number $N$ that names
the dimension of every tensor in this book.}

\section{Why one space will not do}

Every problem in \cref{ch:reductions-pde} reduced to the same abstract shape: find
$u$ in some space such that
\[
  a(u,v) = \ell(v) \qquad \text{for all test functions } v,
\]
with $a$ the bilinear form of \cref{ch:galerkin}. To compute, we replace the
infinite-dimensional space of all admissible fields by a finite-dimensional
subspace $V_h$ spanned by simple, locally-supported basis functions. The unknown
becomes a finite list of coefficients, and $a(u,v)=\ell(v)$ becomes a matrix
equation. This is the finite-element method, and the basis functions are the
\emph{finite elements}.

The temptation is to use one space for everything. After all, an electric field
$\vE$ and a magnetic field $\vB$ are both just vector fields---three numbers at
each point of space. Why not discretize each component the same way we discretize
a scalar, storing values at mesh vertices and interpolating linearly in between?

The answer is that $\vE$ and $\vB$ are not ``just'' vector fields. They are vector
fields with \emph{specific continuity requirements} dictated by Maxwell's
equations, and those requirements differ between the two. Recall from
\cref{ch:maxwell} the interface conditions that hold across any surface separating
two materials: the \emph{tangential} component of $\vE$ is continuous, while the
\emph{normal} component of $\vB$ is continuous. The normal component of $\vE$ may
jump (a surface charge lives in the jump); the tangential component of $\vB$ may
jump (a surface current lives there). A discretization that does not respect the
correct interface condition is not merely inaccurate---it represents a field that
is physically impossible, and the solver will dutifully compute with it.

\begin{keyidea}
The finite-element space must enforce exactly the continuity that physics
requires---no more, no less. Tangentially-continuous fields ($\vE$, the magnetic
vector potential $\vA$) belong in one space; normally-continuous fields (fluxes,
$\vB$, $\vD$) belong in another; scalar potentials belong in a third; quantities
with no inter-element continuity at all belong in a fourth. The degrees of freedom
of each space are chosen to be precisely the quantities that must match across
element faces. \emph{This single principle---DoFs match required continuity---is
the organizing idea of the whole chapter.}
\end{keyidea}

There are four such spaces, and they are not arbitrary. They form a chain---the
\emph{de Rham complex}---linked by the gradient, curl, and divergence operators,
which we develop fully in \cref{ch:derham}. This chapter introduces the four spaces
one at a time, by the degrees of freedom that define them, and shows how the
abstract principle above forces each PDE onto a particular slot in the chain.

\section{Degrees of freedom: the language of finite elements}

Before meeting the four spaces we fix the vocabulary that distinguishes them. A
finite-element space is built from a mesh---a partition of the domain $\dom$ into
simple cells (tetrahedra, hexahedra) sharing faces, edges, and vertices. A field in
the space is determined by a finite list of numbers, its \emph{degrees of freedom}
(DoFs). Each DoF is a \emph{functional}: a rule that reads one number out of a
field.

\begin{definition}[Degree of freedom]
A \emph{degree of freedom} of a finite-element space is a linear functional
$\sigma_i$ on fields, together with a basis function $\phi_j$ satisfying the
\emph{duality} (or \emph{unisolvence}) relation
\[
  \sigma_i(\phi_j) = \delta_{ij}.
\]
A field $u$ in the space is recovered from its DoFs by
$u = \sum_i \sigma_i(u)\,\phi_i$; the coefficient on $\phi_i$ is exactly the value
the functional $\sigma_i$ reads.
\end{definition}

The four spaces differ in \emph{what their functionals read} and \emph{where on the
mesh they live}---and that is the whole story. The choices are summarized in
\cref{tab:fourdofs}, which we then unpack space by space.

\begin{table}[t]
\centering
\caption[Degrees of freedom of the four spaces]{The four finite-element spaces,
distinguished entirely by their degrees of freedom: what each DoF reads and which
geometric entity of the mesh it lives on. The last column anticipates
\cref{ch:derham}: the four spaces form a chain linked by $\Grad$, $\Curl$, $\Div$.}
\label{tab:fourdofs}
\small
\setlength{\tabcolsep}{5pt}
\renewcommand{\arraystretch}{1.3}
\begin{tabular}{@{}lllll@{}}
\toprule
Space & A DoF reads\dots & lives on a\dots & continuity & de Rham slot \\
\midrule
$\Hone$ & a value $u(\vx_i)$ & vertex (node) & full ($C^0$) & start \\
$\Hcurl$ & a circulation $\int_e \mathbf{u}\cdot d\vec\ell$ & edge & tangential & after $\Grad$ \\
$\Hdiv$ & a flux $\int_f \mathbf{u}\cdot\vn\,dA$ & face & normal & after $\Curl$ \\
$\Ltwo$ & a cell value / mean & cell interior & none & after $\Div$ \\
\bottomrule
\end{tabular}
\end{table}

\section{The four spaces}

\subsection{$\Hone$: nodal elements, the home of the scalar potential}

The simplest finite element stores the \emph{value} of a scalar field at each mesh
vertex. The functionals are point evaluations,
\[
  \sigma_i(u) = u(\vx_i), \qquad \vx_i = \text{the $i$-th vertex,}
\]
and the basis functions are the familiar ``hat'' or \emph{Lagrange} functions:
$\phi_i$ is the unique piecewise-polynomial that equals $1$ at vertex $\vx_i$ and
$0$ at every other vertex. A field is reconstructed by interpolating its nodal
values across each cell.

Because neighbouring cells share the vertex $\vx_i$ and both assign it the single
coefficient $\sigma_i(u)=u(\vx_i)$, the reconstructed field takes the \emph{same}
value on both sides of every shared face: it is continuous. This is the space
called $\Hone$, the space of square-integrable scalar fields whose gradient is also
square-integrable. Its discrete subspace is denoted $\Hone_h$, and its conformity is
full continuity ($C^0$): the field has no jumps anywhere.

\begin{physicsbox}
$\Hone$ is the home of the \textbf{electrostatic potential} $V$. The electrostatic
problem of \cref{ch:reductions-pde} solves $-\Div(\varepsilon\Grad V)=0$ for a
scalar $V$ that is continuous everywhere (a potential cannot jump---a jump would be
an infinite field). Nodal continuity is exactly the physics: $V$ must match across
faces, and $\Hone$ enforces precisely that and nothing more.
\end{physicsbox}

The same space serves any scalar field that must be continuous: temperature in heat
conduction, pressure in potential flow, the scalar magnetic potential. Whenever the
unknown is a scalar potential, $\Hone$ is the slot.

\subsection{$\Hcurl$: edge elements, the home of $\vE$ and $\vA$}

Now suppose the unknown is a vector field whose \emph{tangential} component must be
continuous---the electric field $\vE$, or the magnetic vector potential $\vA$. The
naive approach stores all three Cartesian components as separate $\Hone$ fields. We
will see in \cref{sec:spurious} that this is a disaster. The right approach uses a
different kind of DoF entirely.

A \emph{N\'ed\'elec edge element} assigns one DoF to each \emph{edge} of the mesh,
and the functional it reads is the \emph{circulation} of the field along that edge:
\[
  \sigma_e(\mathbf{u}) = \int_e \mathbf{u}\cdot d\vec\ell
  = \int_e \mathbf{u}\cdot\vec t_e \, d\ell,
\]
where $\vec t_e$ is the unit tangent along edge $e$. The basis function $\phi_e$ is
the lowest-order vector field whose circulation is $1$ along edge $e$ and $0$ along
every other edge.

Why does this give tangential continuity? Two cells sharing a face share the edges
bounding that face. Along each shared edge, the circulation DoF is a single number
owned by both cells, so both cells reconstruct a field with the \emph{same
tangential component} along that edge. The tangential component on the shared face
is therefore continuous, while the normal component is free to jump from cell to
cell---exactly the interface behaviour $\vE$ must have. This space is $\Hcurl$, the
space of vector fields whose curl is square-integrable; its discrete subspace
$\Hcurl_h$ is the \emph{edge-element} space.

\begin{physicsbox}
$\Hcurl$ is the home of the \textbf{electric field} $\vE$ and the \textbf{magnetic
vector potential} $\vA$. Every wave and magnetic analysis---eigenmode, driven,
magnetostatic, boundary mode---puts its primary unknown here. The curl--curl
operator $\Curl(\nu\Curl\,\cdot)$ of \cref{ch:reductions-pde} acts on an $\Hcurl$
field, and the tangential-continuity DoFs match the interface condition that
$\vE_{\text{tan}}$ (or $\vA_{\text{tan}}$) be continuous.
\end{physicsbox}

\subsection{$\Hdiv$: face elements, the home of fluxes}

The dual situation is a vector field whose \emph{normal} component must be
continuous---a flux density such as $\vB$, $\vD$, or a current $\vJ$. Here the DoFs
live on \emph{faces}. A \emph{Raviart--Thomas face element} assigns one DoF per face,
reading the \emph{flux} of the field through that face:
\[
  \sigma_f(\mathbf{u}) = \int_f \mathbf{u}\cdot\vn \, dA,
\]
with $\vn$ the unit normal to face $f$. The basis function $\phi_f$ has unit flux
through face $f$ and zero flux through every other face.

Two cells sharing a face share that face's flux DoF---a single number---so both
reconstruct a field with the \emph{same normal component} across the shared face,
while the tangential component may jump. This is $\Hdiv$, the space of vector fields
with square-integrable divergence; its discrete subspace $\Hdiv_h$ is the
\emph{face-element} (Raviart--Thomas) space. Normal continuity is exactly the
condition $\vB$ obeys: $B_n$ continuous, $B_{\text{tan}}$ free.

\subsection{$\Ltwo$: discontinuous elements, the home of densities}

Finally, some quantities have \emph{no} inter-element continuity requirement at all:
a charge density, an energy density, a flux divergence, the right-hand side of a
balance law. For these the appropriate space is $\Ltwo$, merely square-integrable,
with no constraint coupling neighbouring cells. Its discrete subspace $\Ltwo_h$ uses
\emph{discontinuous} elements: the DoFs are cell-interior values or moments, owned by
a single cell and shared with no one. A field in $\Ltwo_h$ may jump freely across
every face.

\subsection{The pattern}

Read \cref{tab:fourdofs} top to bottom and the pattern is unmistakable. As we descend
$\Hone \to \Hcurl \to \Hdiv \to \Ltwo$, the DoFs move from vertices to edges to faces
to cell interiors, and the continuity weakens from full, to tangential, to normal, to
none. \Cref{fig:fourspaces} shows the DoF locations on a single tetrahedron. This
ladder is not a coincidence: the gradient turns a nodal scalar into an edge field, the
curl turns an edge field into a face field, and the divergence turns a face field into
a cell density---the de Rham complex, the subject of \cref{ch:derham}. For now, the
take-away is the placement rule, which we can already use to choose a space for any
field whose continuity we know.

\begin{figure}[t]
\centering
\begin{tikzpicture}[scale=1.0, every node/.style={font=\scriptsize}]
% reusable tetrahedron macro via four corners
\def\A{(0,0)}      \def\B{(2.0,0.25)}   \def\C{(1.0,1.7)}   \def\D{(1.05,0.55)}
% --- panel layout: 2x2 grid of tetrahedra ---
\foreach \px/\py/\ttl/\sub in {%
  0/0/{$\Hone$ (nodal)}/{DoFs on vertices},%
  5/0/{$\Hcurl$ (edge)}/{DoFs on edges},%
  0/-3.4/{$\Hdiv$ (face)}/{DoFs on faces},%
  5/-3.4/{$\Ltwo$ (discontinuous)}/{DoF in interior}%
}{
  \begin{scope}[shift={(\px,\py)}]
    % tetra skeleton (back edges dashed)
    \coordinate (a) at (0,0);
    \coordinate (b) at (2.0,0.25);
    \coordinate (c) at (1.0,1.7);
    \coordinate (d) at (1.05,0.62);
    \draw[simGray,densely dashed] (a)--(d) (b)--(d) (c)--(d);
    \draw[simInk,thick] (a)--(b)--(c)--cycle;
    \node[font=\footnotesize\bfseries, simInk] at (1.0,2.15) {\ttl};
    \node[font=\scriptsize\itshape, simGray] at (1.0,-0.45) {\sub};
  \end{scope}
}
% --- H1: dots on vertices ---
\begin{scope}[shift={(0,0)}]
  \foreach \p in {(0,0),(2.0,0.25),(1.0,1.7),(1.05,0.62)}{
    \fill[simBlue] \p circle (2.6pt);}
\end{scope}
% --- Hcurl: arrows along edges ---
\begin{scope}[shift={(5,0)}]
  \coordinate (a) at (0,0); \coordinate (b) at (2.0,0.25);
  \coordinate (c) at (1.0,1.7); \coordinate (d) at (1.05,0.62);
  \foreach \u/\v in {a/b,b/c,c/a}{
    \draw[-{Stealth[length=1.8mm]},simRust,thick]
      ($(\u)!0.25!(\v)$)--($(\u)!0.75!(\v)$);}
  \foreach \u/\v in {a/d,b/d,c/d}{
    \draw[-{Stealth[length=1.6mm]},simRust,thick,opacity=0.6]
      ($(\u)!0.3!(\v)$)--($(\u)!0.7!(\v)$);}
\end{scope}
% --- Hdiv: a normal arrow through each face (front face shown) ---
\begin{scope}[shift={(0,-3.4)}]
  \coordinate (a) at (0,0); \coordinate (b) at (2.0,0.25);
  \coordinate (c) at (1.0,1.7); \coordinate (d) at (1.05,0.62);
  % front face centroid abc
  \coordinate (fabc) at ($(a)!0.5!(b)$);
  \fill[simGreen!18] (a)--(b)--(c)--cycle;
  \draw[simInk,thick] (a)--(b)--(c)--cycle;
  \node[circle,fill=simGreen,inner sep=1.4pt] (gc) at (1.0,0.65){};
  \draw[-{Stealth[length=2mm]},simGreen,thick] (gc)++(0.0,0.0)--+(0.35,-0.7);
  % a couple of side-face normals
  \node[circle,fill=simGreen,inner sep=1.1pt] (g2) at (0.55,0.55){};
  \draw[-{Stealth[length=1.6mm]},simGreen,thick,opacity=0.6] (g2)--+(-0.45,-0.15);
\end{scope}
% --- L2: single interior dot ---
\begin{scope}[shift={(5,-3.4)}]
  \fill[simViolet] (1.0,0.6) circle (3.0pt);
\end{scope}
\end{tikzpicture}
\caption[The four spaces and their DoF locations]{The hero picture: the same
tetrahedron, four spaces. $\Hone$ places its degrees of freedom on \emph{vertices}
(blue dots: nodal values); $\Hcurl$ on \emph{edges} (rust arrows: tangential
circulations); $\Hdiv$ on \emph{faces} (green normals: fluxes through faces); $\Ltwo$
in the \emph{interior} (violet dot: a single cell value, shared with no neighbour). The
DoF location \emph{is} the continuity: a quantity owned by a shared entity is continuous
across it. Reading the panels in order traces the de Rham ladder of \cref{ch:derham}.}
\label{fig:fourspaces}
\end{figure}

\section{Conformity: why the DoF location is the continuity}

We have asserted that the location of a DoF determines the continuity of the space.
This deserves a precise statement, because it is the mechanism behind everything in
the chapter.

When two cells share a geometric entity---a vertex, an edge, a face---and a DoF lives
\emph{on} that entity, both cells reference the \emph{same single coefficient} for that
DoF. Whatever quantity the functional reads is therefore identical as computed from
either side. The reconstructed fields, built from those coefficients, must agree on
that quantity across the interface. The DoF type fixes \emph{which} quantity:

\begin{itemize}[nosep]
\item a \emph{vertex value} shared $\Rightarrow$ the field value matches: $C^0$
  continuity ($\Hone$);
\item an \emph{edge circulation} shared $\Rightarrow$ the tangential component matches
  along the edge: tangential continuity ($\Hcurl$);
\item a \emph{face flux} shared $\Rightarrow$ the normal component matches across the
  face: normal continuity ($\Hdiv$);
\item \emph{nothing} shared (interior DoF) $\Rightarrow$ no constraint: discontinuity
  ($\Ltwo$).
\end{itemize}

\Cref{fig:tangnormal} draws the two vector cases side by side. The crucial observation
is that $\Hcurl$ and $\Hdiv$ each constrain \emph{exactly one} of the two pieces of a
vector field at an interface (tangential or normal) and leave the other free. This is
not a weakness; it is the point. The interface conditions of Maxwell's equations
constrain exactly one piece and leave the other free, with the freed piece carrying a
physical surface source. A space that constrained \emph{both} pieces (as componentwise
$\Hone$ does) would forbid the surface charges and surface currents that physics
permits---and would over-constrain the solution into the spurious modes of the next
section.

\begin{figure}[t]
\centering
\begin{tikzpicture}[scale=1.0, every node/.style={font=\scriptsize}]
% --- Left: tangential continuity (Hcurl) ---
\begin{scope}
  \node[font=\footnotesize\bfseries] at (1.4,2.2) {$\Hcurl$: tangential continuous};
  \draw[simInk,thick] (1.4,-0.2)--(1.4,1.9);
  \node[simGray] at (1.4,-0.5) {shared face};
  \fill[simBlue!8] (0,-0.2) rectangle (1.4,1.9);
  \fill[simRust!8] (1.4,-0.2) rectangle (2.8,1.9);
  % tangential (vertical) components: equal magnitude both sides
  \draw[-{Stealth[length=2mm]},simBlue,thick] (1.05,0.4)--(1.05,1.1);
  \draw[-{Stealth[length=2mm]},simRust,thick] (1.75,0.4)--(1.75,1.1);
  \node[simBlue] at (0.62,0.75) {$u_{\text{tan}}$};
  \node[simRust] at (2.2,0.75) {$u_{\text{tan}}$};
  % normal components: may differ
  \draw[-{Stealth[length=1.6mm]},simBlue,opacity=0.55] (1.05,1.5)--(0.6,1.5);
  \draw[-{Stealth[length=1.6mm]},simRust,opacity=0.55] (1.75,1.5)--(2.5,1.5);
  \node[simGray] at (1.4,-0.95) {$\;u_{\text{tan}}$ matches, $u_n$ may jump};
\end{scope}
% --- Right: normal continuity (Hdiv) ---
\begin{scope}[shift={(5,0)}]
  \node[font=\footnotesize\bfseries] at (1.4,2.2) {$\Hdiv$: normal continuous};
  \draw[simInk,thick] (1.4,-0.2)--(1.4,1.9);
  \node[simGray] at (1.4,-0.5) {shared face};
  \fill[simBlue!8] (0,-0.2) rectangle (1.4,1.9);
  \fill[simGreen!8] (1.4,-0.2) rectangle (2.8,1.9);
  % normal (horizontal) components: equal both sides
  \draw[-{Stealth[length=2mm]},simBlue,thick] (0.9,0.95)--(1.4,0.95);
  \draw[-{Stealth[length=2mm]},simGreen,thick] (1.4,0.95)--(1.9,0.95);
  \node[simBlue] at (0.55,1.25) {$u_n$};
  \node[simGreen] at (2.25,1.25) {$u_n$};
  % tangential: may differ
  \draw[-{Stealth[length=1.6mm]},simBlue,opacity=0.55] (1.05,0.35)--(1.05,0.0);
  \draw[-{Stealth[length=1.6mm]},simGreen,opacity=0.55] (1.75,0.35)--(1.75,0.75);
  \node[simGray] at (1.4,-0.95) {$\;u_n$ matches, $u_{\text{tan}}$ may jump};
\end{scope}
\end{tikzpicture}
\caption[Tangential versus normal continuity]{The two vector spaces each constrain one
component of a field across a shared face and free the other. \emph{Left:} $\Hcurl$
forces the tangential component $u_{\text{tan}}$ to match (the heavy vertical arrows are
equal) while the normal component $u_n$ may jump. \emph{Right:} $\Hdiv$ forces the
normal component $u_n$ to match (the heavy horizontal arrows are equal) while the
tangential component may jump. These are precisely the Maxwell interface conditions for
$\vE$ (tangential continuous) and $\vB$ (normal continuous) from \cref{ch:maxwell}.}
\label{fig:tangnormal}
\end{figure}

\section{The spurious-mode pitfall: why not just use nodal vectors?}
\label{sec:spurious}

We can now confront the question that motivates edge elements directly. Why not
discretize $\vE$ component by component in $\Hone$---store $E_x$, $E_y$, $E_z$ each as a
nodal scalar field? It is simpler to implement and uses tooling we already have. The
reason it fails is one of the most important lessons in computational electromagnetics.

A nodal vector field forces \emph{all three} components to be continuous at every
vertex, hence the \emph{whole vector} is continuous across faces---both its tangential
and its normal parts. But Maxwell's equations do \emph{not} require the normal part of
$\vE$ to be continuous; across a dielectric interface, or at a re-entrant metallic
corner, $\vE$ has a genuinely discontinuous normal component and a field singularity.
Forcing full continuity where physics demands only tangential continuity over-constrains
the discrete space: the true singular field cannot be represented, and the extra
constraints have to go somewhere.

Where they go is into \emph{spurious modes}---non-physical eigenvectors that pollute the
spectrum of the discrete curl--curl operator. In an eigenmode analysis
(\cref{ch:eigenmodes}) these appear as extra ``resonances'' at frequencies that no
physical cavity possesses, interleaved unpredictably with the real ones, impossible to
filter out reliably. \Cref{fig:spurious} sketches the mechanism at a re-entrant corner:
the nodal-vector field is pinned to a single value at the corner vertex and cannot
follow the field around it, while the edge-element field, owning independent tangential
circulations on the two corner edges, represents the singular field cleanly.

\begin{figure}[t]
\centering
\begin{tikzpicture}[scale=1.0, every node/.style={font=\scriptsize}]
% --- Left: nodal vector at a re-entrant corner ---
\begin{scope}
  \node[font=\footnotesize\bfseries] at (1.3,2.3) {nodal vector (wrong)};
  % L-shaped re-entrant domain
  \fill[simBgGray] (0,0)--(2.4,0)--(2.4,1.0)--(1.2,1.0)--(1.2,2.0)--(0,2.0)--cycle;
  \draw[simInk,thick] (0,0)--(2.4,0)--(2.4,1.0)--(1.2,1.0)--(1.2,2.0)--(0,2.0)--cycle;
  \fill[simRust] (1.2,1.0) circle (2.6pt);
  \node[simRust] at (1.62,1.18) {corner};
  % single pinned vector at corner -> cannot bend
  \draw[-{Stealth[length=2mm]},simInk,thick] (1.2,1.0)--+(0.55,0.25);
  % field tries to wrap but is forced single-valued (X marks)
  \node[simRust,font=\footnotesize] at (0.8,0.55) {$\times$};
  \node[simRust,font=\footnotesize] at (1.85,1.4) {$\times$};
  \node[simGray] at (1.2,-0.5) {spurious: over-constrained};
\end{scope}
% --- Right: edge element at the same corner ---
\begin{scope}[shift={(4.6,0)}]
  \node[font=\footnotesize\bfseries] at (1.3,2.3) {edge element (right)};
  \fill[simBgGreen] (0,0)--(2.4,0)--(2.4,1.0)--(1.2,1.0)--(1.2,2.0)--(0,2.0)--cycle;
  \draw[simInk,thick] (0,0)--(2.4,0)--(2.4,1.0)--(1.2,1.0)--(1.2,2.0)--(0,2.0)--cycle;
  \fill[simGreen] (1.2,1.0) circle (2.6pt);
  % independent edge circulations wrap around the corner
  \draw[-{Stealth[length=1.8mm]},simRust,thick] (1.2,1.0)--+(-0.0,0.55);
  \draw[-{Stealth[length=1.8mm]},simRust,thick] (1.2,1.0)--+(0.6,0.0);
  \draw[simBlue,thick,-{Stealth[length=1.6mm]}]
    (1.7,1.45) arc (35:-120:0.55);
  \node[simBlue] at (2.05,1.0) {$\vE$};
  \node[simGray] at (1.2,-0.5) {singular field captured};
\end{scope}
\end{tikzpicture}
\caption[Spurious modes at a re-entrant corner]{Why componentwise-nodal $\vE$ fails.
\emph{Left:} a nodal-vector field pins the whole vector to one value at the corner
vertex; it forces normal continuity that physics does not require and cannot follow the
field's singular wrap around a re-entrant corner ($\times$ marks the constraints that
have nowhere to go---they become spurious eigenmodes). \emph{Right:} an edge element owns
independent tangential circulations on the corner's edges, leaves the normal component
free, and represents the singular field cleanly. This is the reason every eigenmode and
driven analysis uses $\Hcurl$.}
\label{fig:spurious}
\end{figure}

\begin{pitfall}
Never discretize $\vE$, $\vA$, or any tangentially-continuous field as three
independent nodal ($\Hone$) components. Doing so imposes \emph{normal} continuity that
Maxwell does not require, over-constrains the discrete space, and injects
\textbf{spurious modes}---non-physical eigenvalues that corrupt every eigenmode and
driven solve and cannot be reliably filtered out. Use edge elements ($\Hcurl$), whose
DoFs constrain only the tangential component. This single design decision is why Palace's
eigenmode, driven, magnetostatic, and boundary-mode solvers all assemble their primary
operator over an $\Hcurl$ (\code{ND\_FECollection}) space.
\end{pitfall}

\section{From reference element to global space}

So far we have described DoFs and continuity abstractly. The actual construction has two
layers, which \cref{ch:refelement} develops in full; we sketch them here so the
machinery is in view.

\paragraph{The reference element.} Rather than define basis functions separately on every
cell of the mesh, we define them \emph{once} on a single \emph{reference element}---a
canonical unit tetrahedron or unit cube---and map them to each physical cell by an affine
or isoparametric transformation. Each of the four spaces has its own family of reference
shape functions: nodal Lagrange polynomials for $\Hone$, N\'ed\'elec edge functions for
$\Hcurl$, Raviart--Thomas face functions for $\Hdiv$, and discontinuous Lagrange or
Legendre polynomials for $\Ltwo$. The vector spaces $\Hcurl$ and $\Hdiv$ use special
\emph{covariant} (Piola) and \emph{contravariant} transformations that preserve tangential
and normal components respectively under the mapping---a detail that is what actually makes
the continuity survive the map to a curved or skewed cell.

\paragraph{Gluing into the global space.} The global finite-element space is assembled by
\emph{identifying} the DoFs that shared mesh entities have in common. Each physical edge
contributes one $\Hcurl$ DoF no matter how many cells border it; each physical face
contributes one $\Hdiv$ DoF; each vertex one $\Hone$ DoF. The local-to-global \emph{DoF
map}---a table recording, for each cell, which global DoF each of its local DoFs maps
to---is what stitches the per-cell reference bases into one continuous (or
tangentially/normally continuous) global field. In the matrix-free pipeline of
\cref{ch:matrixfree} this map is the \emph{element restriction} operator $\mathsf{G}$ that
gathers a global DoF vector into per-element local-dof tensors and scatters the result
back; the per-element local-dof axis is the $L$ axis of the element-local tensor
introduced there.

In Palace, this whole construction is a single object. The function space is built by
pairing a mesh with a \emph{collection} that names the family and order:
\begin{center}
\code{fe\_space(mesh, collection)} $\longrightarrow$
\code{FiniteElementSpace[N]},
\end{center}
where the collection's type ($\Hone$, $\Hcurl$, $\Hdiv$, or $\Ltwo$) selects the de Rham
slot and $N$ is the resulting dimension---the subject of \cref{sec:truedof}. The mesh and
the collection are independent inputs: the same collection over a finer mesh, or the same
mesh under a higher-order collection, are both well-formed spaces. These two axes of
refinement are the topic of the next section.

\section{Polynomial order and the \code{FECollection} family}

Each space comes in a tower of \emph{orders}. The order-$p$ version of a space replaces the
lowest-order shape functions with degree-$p$ polynomials, adding more DoFs per cell and
representing the field more accurately. Palace organizes the four families and their orders
through \emph{finite-element collections}: \code{H1\_FECollection}, \code{ND\_FECollection}
(N\'ed\'elec, $\Hcurl$), \code{RT\_FECollection} (Raviart--Thomas, $\Hdiv$), and
\code{L2\_FECollection}. A collection is the object passed as the second argument to
\code{fe\_space} above; its \emph{type} fixes the de Rham slot and its order fixes the
polynomial degree.

\Cref{tab:collections} collects the four families with the data that matters for building a
solver: the de Rham slot, the geometric entity that carries the DoFs, and the
\emph{minimum order} $p_{\min}$ each family admits.

\begin{table}[t]
\centering
\caption[The four \code{FECollection} families]{The four finite-element collection
families. The minimum order $p_{\min}$ differs by family: $\Hone$ and $\Hcurl$ require
$p_{\min}=1$ (a constant nodal scalar or a lowest-order edge function is the floor), while
$\Hdiv$ and $\Ltwo$ admit $p_{\min}=0$ (a piecewise-constant flux or cell value is
meaningful). These floors bound the coarsening sequence of the $p$-multigrid hierarchy.}
\label{tab:collections}
\small
\setlength{\tabcolsep}{6pt}
\renewcommand{\arraystretch}{1.3}
\begin{tabular}{@{}lllcl@{}}
\toprule
Space & Collection & DoF entity & $p_{\min}$ & primary field \\
\midrule
$\Hone$  & \code{H1\_FECollection} & vertices (nodes) & $1$ & scalar potential $V$ \\
$\Hcurl$ & \code{ND\_FECollection} & edges            & $1$ & $\vE$, $\vA$ \\
$\Hdiv$  & \code{RT\_FECollection} & faces            & $0$ & flux $\vB$, $\vD$ \\
$\Ltwo$  & \code{L2\_FECollection} & cell interiors   & $0$ & density, $B=\Curl\vE$ (2-D) \\
\bottomrule
\end{tabular}
\end{table}

\paragraph{$h$- versus $p$-refinement.} There are two independent ways to make a finite-element
solution more accurate, drawn in \cref{fig:hp}. \emph{$h$-refinement} subdivides the mesh into
smaller cells (the mesh size $h$ shrinks) while keeping the polynomial order fixed; it adds
resolution where the field varies rapidly and is the basis of the adaptive mesh refinement of
\cref{ch:amr}. \emph{$p$-refinement} raises the polynomial order $p$ on a fixed mesh, adding DoFs
per cell; on smooth fields it converges far faster (exponentially, rather than algebraically) than
$h$-refinement. The two are orthogonal---one varies the mesh argument of \code{fe\_space}, the
other the collection argument---and they can be combined ($hp$-refinement).

\begin{figure}[t]
\centering
\begin{tikzpicture}[scale=1.0, every node/.style={font=\scriptsize}]
% baseline single coarse element
\begin{scope}
  \node[font=\footnotesize\bfseries] at (1.0,1.65) {coarse, $p=1$};
  \draw[simInk,thick] (0,0) rectangle (2,1.2);
  \foreach \x/\y in {0/0,2/0,2/1.2,0/1.2}{\fill[simBlue] (\x,\y) circle (2pt);}
\end{scope}
% h-refinement: subdivide
\begin{scope}[shift={(3.4,0)}]
  \node[font=\footnotesize\bfseries] at (1.0,1.65) {$h$-refine ($p=1$)};
  \draw[simInk,thick] (0,0) rectangle (2,1.2);
  \draw[simGray] (1,0)--(1,1.2) (0,0.6)--(2,0.6);
  \foreach \x in {0,1,2}{\foreach \y in {0,0.6,1.2}{\fill[simBlue] (\x,\y) circle (2pt);}}
  \node[simGray] at (1.0,-0.5) {smaller cells};
\end{scope}
% p-refinement: more DoFs, same cell
\begin{scope}[shift={(6.8,0)}]
  \node[font=\footnotesize\bfseries] at (1.0,1.65) {$p$-refine ($p=2$)};
  \draw[simInk,thick] (0,0) rectangle (2,1.2);
  \foreach \x in {0,1,2}{\foreach \y in {0,0.6,1.2}{\fill[simRust] (\x,\y) circle (2pt);}}
  \node[simGray] at (1.0,-0.5) {higher order};
\end{scope}
\end{tikzpicture}
\caption[$h$- versus $p$-refinement]{Two orthogonal routes to accuracy, shown for a nodal
($\Hone$) element. \emph{$h$-refinement} (middle) subdivides the cell, shrinking the mesh size $h$
while keeping order $p=1$; the DoFs (blue) multiply by subdivision. \emph{$p$-refinement} (right)
keeps the single cell but raises the order to $p=2$, adding interior and edge DoFs (rust). Palace's
geometric $p$-multigrid (\cref{ch:multigrid}) builds a hierarchy of order-coarsened collections; the
$p_{\min}$ floors of \cref{tab:collections} bound how far the order can drop.}
\label{fig:hp}
\end{figure}

\paragraph{The $p$-multigrid order schedule.} Palace's geometric multigrid preconditioner
(\cref{ch:multigrid}) needs not a single space but a \emph{tower} of them, one per multigrid level,
differing in order. The collection-schedule operator produces a list of collections from the finest
order $p$ down to the floor $p_{\min}$, coarsening the order level by level. Two coarsening policies
are offered: \emph{linear}, which steps $p \mapsto p-1$ each level, and \emph{logarithmic} (the
default), which halves the gap to the floor, $p \mapsto \lfloor (p+p_{\min})/2 \rfloor$, reaching the
floor in far fewer levels for large $p$. The schedule terminates at $p_{\min}$---which is why the
floor differs by family: an $\Hone$ or $\Hcurl$ tower bottoms out at order $1$ (a constant scalar or
lowest edge element), while an $\Hdiv$ or $\Ltwo$ tower can descend to the piecewise-constant order
$0$. When the tower has a single level, the schedule degenerates to one ordinary \code{fe\_space}
construction.

\section{True degrees of freedom: the number $N$}
\label{sec:truedof}

We close with the single most-used number in this book. The dimension of the discrete space
$V_h$---the number of basis functions, hence the length of the coefficient vector, hence the size
of every assembled matrix---is denoted
\[
  N = \dim V_h.
\]
This is the $N$ in every \code{Tensor[N]} and every \code{LinearOperator[N,N]} throughout the
framework chapters; the function-space construction is the operator that \emph{defines} $N$.

There is a subtlety. The raw count of DoFs over all cells---summing local DoFs cell by cell---over-counts, because DoFs on shared entities are identified. Worse, on an adaptively refined mesh, a
\emph{hanging node} (a vertex of a fine cell lying in the middle of a coarse cell's face) is not an
independent DoF: its value is \emph{constrained} to be an interpolation of the surrounding coarse
DoFs, so that the global field stays conforming. The space therefore distinguishes two counts:

\begin{itemize}[nosep]
\item the \emph{local} (or \emph{L-vector}) DoF count, summed over cells with shared and hanging
  DoFs still counted multiply---the natural size for per-element work;
\item the \emph{true} DoF count $N$, with shared DoFs identified once and hanging-node constraints
  eliminated---the number of genuinely independent coefficients, the dimension of $V_h$.
\end{itemize}

It is the \emph{true}-DoF count $N$ that the solver works with: the assembled operators are $N\times
N$, the Krylov vectors of \cref{ch:krylov} are length $N$, the unknowns are $N$ true coefficients.
The transfer between the two pictures---gather the $N$ true coefficients into the larger local
vector, apply the constraints, scatter back---is exactly the prolongation/restriction pair the space
carries, and exactly the boundary the element-restriction operator $\mathsf{G}$ of
\cref{ch:matrixfree} sits on. For the un-refined, single-order meshes of most of this book there are
no hanging nodes, and $N$ is simply ``DoFs with shared entities counted once.''

\begin{notationbox}
$N = \dim V_h$ is the \textbf{true-DoF count}: the number of independent coefficients after shared
DoFs are identified and hanging-node constraints are removed. It is the dimension of every
\code{Tensor[N]} and the size of every \code{LinearOperator[N,N]} in the book. Different analyses on
the same mesh have different $N$, because they use different spaces---an $\Hone$ scalar problem and
an $\Hcurl$ vector problem over one mesh produce entirely different $N$.
\end{notationbox}

\section{Worked examples}

\subsection{Counting degrees of freedom}

The placement rule of \cref{tab:fourdofs} makes the lowest-order DoF count a purely combinatorial
exercise: count vertices for $\Hone$, edges for $\Hcurl$, faces for $\Hdiv$, cells for $\Ltwo$. We do
it first on a single tetrahedron, then on a small two-element mesh, to see how sharing changes the
arithmetic.

\begin{workedexample}{Lowest-order DoFs on a tetrahedron and a 2-element mesh}{dofcount}
A single tetrahedron has $4$ vertices, $6$ edges, $4$ faces, and $1$ cell. For the lowest-order
spaces the DoF count is the count of the carrying entity:
\[
  N_{\Hone} = 4, \quad N_{\Hcurl} = 6, \quad N_{\Hdiv} = 4, \quad N_{\Ltwo} = 1.
\]
So even on one cell the three first spaces have genuinely different sizes---vertices, edges, and faces
of a tetrahedron number $4$, $6$, $4$. An $\Hcurl$ problem is the largest; this is the price of getting
the continuity right.

Now glue a \emph{second} tetrahedron onto the first across a shared triangular face. The two cells share
that face's $3$ vertices and $3$ edges and the face itself; everything else is new. Counting the
combined mesh:
\begin{itemize}[nosep]
\item \textbf{vertices:} $4 + 1 = 5$. (The second tet adds one new apex; its other three vertices are
  the shared face.) So $N_{\Hone} = 5$.
\item \textbf{edges:} the first tet has $6$; the second adds the $3$ edges from its new apex to the
  shared face's vertices, the other $3$ being shared. So $6 + 3 = 9$, giving $N_{\Hcurl} = 9$.
\item \textbf{faces:} the first tet has $4$; the second adds $3$ new faces (the fourth is the shared
  face). So $4 + 3 = 7$, but the \emph{shared} face is counted once: $N_{\Hdiv} = 7$.
\item \textbf{cells:} $2$, so $N_{\Ltwo} = 2$.
\end{itemize}
The naive ``sum over cells'' count would give $\Hone$: $4+4=8$, $\Hcurl$: $6+6=12$, $\Hdiv$:
$4+4=8$. The true counts $5,9,7$ are smaller because the shared vertices, edges, and the shared face
are \emph{identified}---each contributes a single global DoF. That identification is exactly the
gluing of \cref{sec:truedof} and the source of the gap between the local and true DoF counts. The
amount of sharing differs by space: the shared face removes $3$ from the vertex count, $3$ from the
edge count, but only $1$ from the face count.
\end{workedexample}

\subsection{A lowest-order edge basis function}

The defining feature of $\Hcurl$---tangential support on a single edge---can be checked by hand on the
reference triangle. We exhibit the lowest-order N\'ed\'elec basis function for one edge and verify its
tangential trace.

\begin{workedexample}{The N\'ed\'elec edge function and its tangential trace}{whitney}
Take the reference triangle with vertices $\vx_0=(0,0)$, $\vx_1=(1,0)$, $\vx_2=(0,1)$, and let
$\lambda_0,\lambda_1,\lambda_2$ be its barycentric (nodal $\Hone$) coordinates---the hat functions with
$\lambda_i(\vx_j)=\delta_{ij}$. Explicitly $\lambda_0 = 1-x-y$, $\lambda_1 = x$, $\lambda_2 = y$, and
note $\Grad\lambda_0=(-1,-1)$, $\Grad\lambda_1=(1,0)$, $\Grad\lambda_2=(0,1)$.

The lowest-order edge (Whitney) function associated with the edge $e_{01}$ joining $\vx_0$ and $\vx_1$
is
\[
  \boldsymbol{\phi}_{01} \;=\; \lambda_0\,\Grad\lambda_1 \;-\; \lambda_1\,\Grad\lambda_0
  \;=\; (1-x-y)\,(1,0) \;-\; x\,(-1,-1)
  \;=\; \bigl(\,1-y,\; x\,\bigr).
\]
This is a genuine vector field on the triangle. We check the three edge circulations
$\sigma_e(\boldsymbol{\phi}_{01})=\int_e \boldsymbol{\phi}_{01}\cdot d\vec\ell$.

\smallskip\noindent\emph{Edge $e_{01}$} (its own edge), from $\vx_0=(0,0)$ to $\vx_1=(1,0)$: here
$y=0$, the tangent is $d\vec\ell=(dx,0)$, and $\boldsymbol{\phi}_{01}=(1-0,\,x)=(1,x)$, so
$\boldsymbol{\phi}_{01}\cdot d\vec\ell = 1\,dx$. The circulation is
$\int_0^1 1\,dx = 1.$

\smallskip\noindent\emph{Edge $e_{02}$}, from $\vx_0=(0,0)$ to $\vx_2=(0,1)$: here $x=0$, the tangent is
$(0,dy)$, and $\boldsymbol{\phi}_{01}=(1-y,\,0)$, so the dot product
$\boldsymbol{\phi}_{01}\cdot d\vec\ell = 0$. The circulation is $0.$

\smallskip\noindent\emph{Edge $e_{12}$}, from $\vx_1=(1,0)$ to $\vx_2=(0,1)$: parametrize
$(x,y)=(1-t,\,t)$ for $t\in[0,1]$, so $d\vec\ell=(-1,1)\,dt$ and
$\boldsymbol{\phi}_{01}=(1-t,\,1-t)$. Then
$\boldsymbol{\phi}_{01}\cdot d\vec\ell = (1-t)(-1)+(1-t)(1) = 0$,
and the circulation is $0.$

\smallskip So $\sigma_{e_{01}}(\boldsymbol{\phi}_{01})=1$ and
$\sigma_{e_{02}}=\sigma_{e_{12}}=0$: the function satisfies the duality relation
$\sigma_e(\boldsymbol{\phi}_{01})=\delta_{e,e_{01}}$. Its tangential trace is supported on \emph{exactly
one} edge. Consequently, when two triangles share edge $e_{01}$, both reference $\boldsymbol{\phi}_{01}$
through the \emph{same} circulation coefficient, and the assembled field has continuous tangential
component along $e_{01}$ and nowhere else constrained---tangential continuity, edge by edge, exactly as
\cref{fig:tangnormal} promised.
\end{workedexample}

\summarysection
A finite-element discretization must place each field in a space whose degrees of freedom enforce
exactly the continuity the physics requires. There are four canonical spaces, distinguished entirely by
their DoFs: $\Hone$ stores nodal \emph{values} on vertices and is fully continuous (the home of the
scalar potential $V$); $\Hcurl$ (N\'ed\'elec edge elements) stores \emph{circulations} on edges and is
tangentially continuous (the home of $\vE$ and $\vA$); $\Hdiv$ (Raviart--Thomas face elements) stores
\emph{fluxes} on faces and is normally continuous (the home of $\vB$, $\vD$); $\Ltwo$ stores
\emph{cell} quantities and has no inter-element continuity (the home of densities). Because a DoF on a
shared entity is a single coefficient owned by both neighbours, the DoF \emph{location} is the
continuity. Discretizing $\vE$ as componentwise nodal vectors over-constrains it---it imposes normal
continuity physics does not require---and injects spurious modes that ruin eigenmode and driven solves;
edge elements get the continuity exactly right, which is why every magnetic and wave analysis uses
$\Hcurl$. Each space comes in a tower of polynomial orders with family-dependent floors
($p_{\min}=1$ for $\Hone$/$\Hcurl$, $0$ for $\Hdiv$/$\Ltwo$), assembled from reference shape functions
glued by a local-to-global DoF map; refinement is $h$ (smaller cells) or $p$ (higher order). The
dimension $N=\dim V_h$---the true-DoF count, with shared DoFs identified and hanging-node constraints
removed---is the $N$ that names every tensor in the book.

\furtherreading
The definitive reference for finite-element spaces in electromagnetics is
Monk~\citep{monk2003}, which develops $\Hcurl$ and $\Hdiv$, the N\'ed\'elec and Raviart--Thomas
elements, and the spurious-mode question with full rigour. The unifying structural view---the four
spaces as a discrete de Rham complex with exactness---is the subject of finite-element exterior
calculus; Arnold, Falk, and Winther~\citep{arnold2006feec} is the canonical account and the natural
sequel to \cref{ch:derham}. For an engineering-oriented treatment tied directly to computational
electromagnetics, with the spurious-mode pathology worked through on real cavity problems,
Jin~\citep{jin2014} is the standard text. The continuity arguments of \cref{sec:spurious} are made
precise there; the tangential/normal trace theory underlying them is Monk's Chapter~3.

\exercisessection
\begin{enumerate}[nosep]
\item For each of the four spaces, state in one sentence which Maxwell interface condition its DoFs
  enforce, and name a physical field that belongs in it. (Use \cref{tab:fourdofs} and the
  \texttt{physicsbox} anchors.)
\item A single hexahedron (cube) has $8$ vertices, $12$ edges, $6$ faces, and $1$ cell. Give the
  lowest-order DoF count $N$ for each of the four spaces on one hexahedron. Which space is largest?
\item Two hexahedra are glued across a shared square face. Following
  \cref{wex:dofcount}, compute the true DoF count $N$ for $\Hone$, $\Hcurl$, $\Hdiv$, and $\Ltwo$ on the
  combined mesh, and compare each to the naive sum-over-cells count. How many DoFs does the shared face
  remove from each space?
\item Explain, in terms of degrees of freedom, why discretizing $\vE$ as three independent $\Hone$
  scalar fields imposes \emph{more} continuity than $\Hcurl$ does. Which interface condition does the
  nodal-vector approach wrongly enforce, and what physical surface quantity does that forbid?
\item The $p$-multigrid order schedule for an $\Hcurl$ ($p_{\min}=1$) space at finest order $p=8$
  produces a tower of collections. List the order sequence under (a) linear coarsening
  ($p\mapsto p-1$) and (b) logarithmic coarsening ($p\mapsto\lfloor(p+1)/2\rfloor$). How many levels
  does each policy produce before reaching the floor?
\item Repeat the previous exercise for an $\Hdiv$ ($p_{\min}=0$) space at $p=8$. Why can this tower
  descend one order further than the $\Hcurl$ tower, and what does the order-$0$ space represent
  physically?
\item Verify, by the method of \cref{wex:whitney}, that the edge function
  $\boldsymbol{\phi}_{02}=\lambda_0\Grad\lambda_2-\lambda_2\Grad\lambda_0$ on the reference triangle has
  circulation $1$ along edge $e_{02}$ and $0$ along $e_{01}$ and $e_{12}$.
\item The number $N=\dim V_h$ differs between an electrostatic ($\Hone$) and an eigenmode ($\Hcurl$)
  analysis run on the \emph{same} mesh. Using the single-tetrahedron and two-tet counts of
  \cref{wex:dofcount}, explain why, and state which analysis produces the larger linear system on a given
  mesh.
\end{enumerate}
